% Generated draft: requires mathematical and source review.
\documentclass[11pt]{article}
\usepackage[margin=25mm]{geometry}
\usepackage{fontspec}
\setmainfont{Latin Modern Roman}
\usepackage{amsmath,amssymb,mathtools,unicode-math}
\setmathfont{Latin Modern Math}
\usepackage{enumitem,needspace,xurl,hyperref}
\setlength{\parindent}{0pt}
\setlength{\parskip}{5pt}
\setlength{\emergencystretch}{3em}
\newcommand{\N}{\mathbb{N}}
\newcommand{\Z}{\mathbb{Z}}
\newcommand{\Q}{\mathbb{Q}}
\newcommand{\R}{\mathbb{R}}
\newcommand{\C}{\mathbb{C}}
\newcommand{\F}{\mathbb{F}}
\newcommand{\A}{\mathbb{A}}
\newcommand{\PP}{\mathbb{P}}
\newcommand{\E}{\mathbb{E}}
\newcommand{\eps}{\varepsilon}
\newcommand{\dd}{\,\mathrm{d}}
\begin{document}
\section*{UnsolvedMath ID: 20000446}
\textbf{Source problem code:} catalogue entry
\subsection{Definitions and mathematical statement}
For a smooth projective surface \(X/\F_2\), write \(\rho_{\mathrm{geom}}(X)=\operatorname{rank}\operatorname{NS}(X_{\overline{\F}_2})\). The arithmetic Picard number concerns classes over \(\F_2\), or equivalently their invariant rational N\'eron--Severi space; it need not equal the geometric number in general. Voloch's question is to exhibit a homogeneous quintic \(f\in\F_2[x,y,z,w]\) such that \(X=V(f)\subset\PP^3_{\F_2}\) is smooth and has Picard number one. The supplied solution uses the stronger geometric rank-one condition, which also gives arithmetic rank one because the hyperplane class is defined over \(\F_2\).

The source's later report attributes to Abbott, Kedlaya and Roe the explicit polynomial
\[
\begin{aligned}
 f={}&x^5+x^3yz+x^2y^2w+x^2yz^2+x^2z^3+xz^2w^2\\
    &+y^5+y^3zw+y^2zw^2+yzw^3+z^5+z^2w^3+w^5.
\end{aligned}
\]
It reports that \(V(f)\) is smooth and \(\rho_{\mathrm{geom}}(V(f))=1\), with the Picard conclusion supplied by their certified finite-precision Frobenius-corank computation. The report labels this a known solution, not a newly proved Picard-rank computation in the present draft.

For its separate candidate smoothness certificate let \(R=\F_2[x,y,z,w]\), let \(R_m\) be the space of homogeneous degree-\(m\) polynomials, and define the Jacobian Macaulay map
\[
 M_m:R_{m-4}^{\,4}\longrightarrow R_m,\quad
 (a_x,a_y,a_z,a_w)\longmapsto
 a_xf_x+a_yf_y+a_zf_z+a_wf_w.
\]
The supplied exact elimination results for this \(f\) are
\[
 \operatorname{rank}M_{13}=560\quad(560\text{-by-}880\text{ matrix}),\qquad
 \operatorname{rank}M_{12}=454\quad(455\text{-by-}660\text{ matrix}).
\]
The candidate proposition further asserts that degree thirteen is the sharp uniform Jacobian-surjectivity threshold for smooth quintic surfaces in \(\PP^3\) in characteristic not dividing five. Its cited complete-intersection Hilbert series is
\[
 \frac{(1-t^4)^4}{(1-t)^4}=(1+t+t^2+t^3)^4,
\]
with top degree twelve. The matrices, elimination certificate and full proof are absent from the supplied record, so these are reported verification data requiring review.

The historical remarks connect the example to error-correcting codes, state second Betti number \(53\), and describe zeta-function point counts to roughly \(\F_{2^{28}}\) as costly. Their reference to ``Tate's theorem'' supplies no precise theorem or hypotheses, and is not used here to assert an unconditional equality of analytic and algebraic ranks. Raskind mentions Shioda's characteristic-zero rank-one examples as possible sources; Voloch says this had been tried once. These are historical suggestions, not a general preservation claim under reduction.
\subsection{Short English statement}
Find a smooth quintic surface over \(\F_2\) with Picard number one. The supplied record reports a known explicit Abbott--Kedlaya--Roe example with geometric Picard number one, plus a separate exact Jacobian-matrix smoothness certificate. These source claims remain subject to mathematical review of the cited computation.
\subsection{Sources}
\begin{itemize}
\item Local UnsolvedMath registry, permanent ID \(20000446\), source code \texttt{catalogue entry}; AIM algebraic number theory corpus, \texttt{aim-algebraic-number-theory-notes.json}, notes[97]. Attribution: Ulam AI / UnsolvedMath Contributors, CC BY 4.0.
\item AIM \emph{Rational and integral points on higher dimensional varieties}, Question 13, attributed to Voloch, with remarks by Voloch and Raskind: \url{https://aimath.org/WWN/qptsurface2/qptsurface2.pdf}. The supplied verification says the official HTML/PDF resolve the damaged finite-field subscript; no new inspection is claimed. The remarks mention Shioda and end with ``Problem/''.
\item The report attributes the explicit quintic and certified Frobenius-corank Picard-rank proof to Abbott, Kedlaya and Roe (AKR), but supplies no full bibliographic citation or computation artifact. The original literature field is empty.
\item Registry dated August 14, 2026: statement exact, status solved, classification \texttt{SOLVED-IN-LITERATURE}; attempt 1 labelled full solution, problem status solved, rigor proved, novelty confidence low. The accompanying research report and candidate exact-elimination contribution are machine-generated aids, not peer reviewed, and require independent verification before citation.
\end{itemize}
\end{document}
